%------------------------------------------------------------------------------%
\begin{exercises}

% Thomas 10.7 pg 52 41-48
%------------------------------------------------------------------------------%
\begin{exercise}
	Utilize o Teorema~\ref{teo:Mudanca-variaveis-serie-potencias}
	para encontrar o intervalo de convergência da série. Encontre também
	a soma da série, como uma função de $x$, dentro do intervalo de convergência.
	\begin{mathlist}
		\item \sum_{n=0}^{\infty} 3^n x^n
		\item \sum_{n=0}^{\infty} \big( e^x-4 \big)^n
		\item \sum_{n=0}^{\infty} \frac{(x-1)^{2n}}{4^n}
		\item \sum_{n=0}^{\infty} \frac{(x+1)^{2n}}{9^n}
		\item \sum_{n=0}^{\infty} \left(\frac{\sqrt{x}}{2}-1\right)^n
		\item \sum_{n=0}^{\infty} \big( \ln x \big)^n
		\item \sum_{n=0}^{\infty} \left(\frac{x^2+1}{3}\right)^n
		\item \sum_{n=0}^{\infty} \left(\frac{x^2-1}{2}\right)^n
	\end{mathlist}
\end{exercise}

%------------------------------------------------------------------------------%
%\begin{exercise}
%\begin{answer}
%\end{answer}
%\end{exercise}

%------------------------------------------------------------------------------%
%\begin{exercise}
%\begin{mathlist}[2]
%  \item
%  \item
%  \item
%\end{mathlist}
%\begin{answer}
%  \begin{alphalist}[3]
%    \item
%    \item
%    \item
%  \end{alphalist}
%\end{answer}
%\end{exercise}

\end{exercises}
%------------------------------------------------------------------------------%
